\section{A gallery of groups}

Definition~\ref{def:group} is short, so the way to make it mean something is to run it against
examples until the pattern is automatic. For each example below the routine is the same: say what
the \emph{set} is, say what the \emph{rule} is, check closure and the three axioms, and only then
collect the properties that make the example worth knowing.

\subsection{Dihedral groups: symmetries of a polygon}\label{sec:dihedral}

\emph{The set.} All rigid motions of a regular $n$-gon in the plane that carry it onto itself. There
are $2n$ of them: $n$ rotations, by multiples of $2\pi/n$ including the identity, and $n$
reflections. The group is written $D_n$.

\emph{The rule.} Do one, then the other, exactly as for the cube.

Reflections are allowed here, and it is worth being clear why this is not inconsistent with
Section~\ref{sec:def}. A flat polygon lives in the plane but sits in space, so you can pick it up,
turn it over and put it back; the motion is something your hands can do. For the solid cube there
is no such trick.

\subsubsection*{Closure}

Two reflections in a row are not a reflection, but they are certainly a symmetry, so nothing leaves
the set. The same picture makes a second point that will matter later: the order of the two
reflections changes the answer.

\sfigs{0.86}{Closure in $D_5$, and the failure of commutativity. Reflecting in the axis of $t$ and then in
the axis of $s$ gives one click of rotation anticlockwise; doing the two reflections in the other
order gives one click clockwise. So $st\neq ts$.}{\figXq}

\subsubsection*{Identity and inverses}

``Leave the polygon alone'' is the rotation by $0^\circ$. A rotation is undone by turning back, which
is the same as turning forward the rest of the way, so the inverse of one click is four clicks when
$n=5$. A reflection undoes itself: $t^2=1$, hence $t\inv=t$.

\sfigs{0.86}{Identity in $D_5$: start, after the reflection $t$, then after doing nothing.}{\figXr}

\sfigs{0.86}{Inverses in $D_5$: one click of rotation, then four more clicks, returns every corner to its
starting position, so $(st)\inv=(st)^4$.}{\figXs}

\subsubsection*{Associativity}

Free of charge, by the remark on page~\pageref{def:group}: the elements are functions on the polygon
and the rule is composition. So a word like $tst$ has exactly one meaning.

\subsubsection*{The extra properties}

$|D_n|=2n$, and $D_n$ is not abelian for $n\ge3$. The two reflections $s$ and $t$ drawn in the
figures \defn{generate} the group, meaning every element is a product of copies of $s$ and $t$.

\sfig{\label{fig:pentaxes}The two reflection axes of $s$ and $t$ for the pentagon. They meet at the
angle $\pi/n$, which is $36^\circ$ when $n=5$, and that is what makes the product $st$ a single
click of rotation.}{\figpent}

\begin{why}
\emph{Why $st$ is one click.} Reflecting in two lines that meet at angle $\theta$ is the same as
rotating by $2\theta$ about the point where they cross. The axes of $s$ and $t$ meet at $\pi/n$, so
$st$ is the rotation by $2\pi/n$: one click. Consequently $(st)^m$ is $m$ clicks, which gives all
$n$ rotations as $m$ runs from $0$ to $n-1$, and $(st)^n=1$.
\end{why}

\begin{why}
\emph{Why every reflection is $s(st)^m$.} A reflection followed by a rotation is again a reflection,
in a different axis; and since $s\cdot s=1$, the word $s(st)^m$ can be rewritten as $t(st)^{m-1}$.
Running $m$ through $0,\dots,n-1$ produces $n$ distinct reflections, whose axes are $\pi/n$ apart.
Together with the rotations this is $2n$ elements, which is all of them.
\end{why}

Here is the whole group $D_5$ written out, both as pictures and as words.

\sfig{The ten symmetries of a regular pentagon. Top row: the five rotations, $(st)^m$ being $m$
clicks of $72^\circ$. Bottom row: the five reflections, with their axes dashed. Each circled number
shows where that corner has been sent.}{\figXv}

\begin{center}
\small
\begin{tabular}{@{}llll@{}}
\toprule
word in $s,t$ & type & what it does & corners $1\,2\,3\,4\,5$ go to\\
\midrule
$1$ & rotation & by $0^\circ$ (0 clicks) & 1 2 3 4 5\\
$st$ & rotation & by $72^\circ$ (1 click) & 2 3 4 5 1\\
$(st)^2$ & rotation & by $144^\circ$ (2 clicks) & 3 4 5 1 2\\
$(st)^3$ & rotation & by $216^\circ$ (3 clicks) & 4 5 1 2 3\\
$(st)^4$ & rotation & by $288^\circ$ (4 clicks) & 5 1 2 3 4\\
\midrule
$s$ & reflection & axis at $126^\circ$ from horizontal & 2 1 5 4 3\\
$t$ & reflection & axis at $90^\circ$ from horizontal & 1 5 4 3 2\\
$tst$ & reflection & axis at $54^\circ$ from horizontal & 5 4 3 2 1\\
$tstst$ & reflection & axis at $18^\circ$ from horizontal & 4 3 2 1 5\\
$tststst$ & reflection & axis at $-18^\circ$ from horizontal & 3 2 1 5 4\\
\bottomrule
\end{tabular}
\end{center}

\begin{pitfall}
``$s$ and $t$ generate $D_n$'' does not mean every element \emph{is} $s$ or $t$. It means every
element is a product of copies of them. Inverses need not be mentioned here because $s\inv=s$ and
$t\inv=t$, but in general a generating set is allowed to use inverses as well.
\end{pitfall}

The rotations of $D_n$, taken alone, are closed under the rule and form a group in their own right:
a \defn{subgroup} of $D_n$, of size $n$. The reflections alone are not, since the product of two of
them is a rotation. Finally, the three facts $s^2=1$, $t^2=1$ and $(st)^n=1$ turn out to determine
$D_n$ completely, in a sense made precise in Section~\ref{sec:presentations}.

\sfig{The eight symmetries of a square, which is $D_4$. Corners are numbered $1,2,3,4$ and each
picture shows where the numbers end up; $r$ is a quarter turn and $s$ a reflection, and the other
three reflections are $sr$, $sr^2$ and $sr^3$.}{\figsquares}

\sfig{The smallest interesting case, $D_3$: three rotations and three reflections of an equilateral
triangle, six symmetries in all.}{\figtriangle}

\begin{tryit}
Multiply two rows of the table by composing their corner maps, right-hand factor first. For example
$t\cdot s$: first $s$ sends $1\,2\,3\,4\,5$ to $2\,1\,5\,4\,3$, then $t$ acts on the result. You
should land on $(st)^4$, the rotation by $-72^\circ$, confirming $ts=(st)\inv$. Then find which word
is the reflection whose axis passes through corner $3$.
\end{tryit}

\subsection{Symmetric groups: rearranging $n$ things}\label{sec:symmetric}

\emph{The set.} All bijections from $\{1,2,\dots,n\}$ to itself. Such a bijection is called a
\defn{permutation}, and the set of them is written $S_n$.

\emph{The rule.} Composition of functions, with the convention of Definition~\ref{def:order}: in
$\tau\sigma$, the permutation $\sigma$ acts first.

Nothing geometric has to be respected here, which is the point of the example. If you think of
$\{1,\dots,n\}$ as $n$ loose points with no distances and no shape, then every rearrangement of them
is a symmetry, and $S_n$ is the symmetry group of a set of $n$ points. In the pictures below the
points are drawn as $n$ fixed \emph{slots}, and a coloured disc is the token currently in that slot.

\subsubsection*{Closure, identity, inverses, associativity}

A rearrangement followed by a rearrangement is a rearrangement, since every slot still ends up with
exactly one token. The identity is the permutation that fixes every slot, written $1$ or $()$. The
inverse of a permutation is obtained by sending every token back where it came from, which on the
picture means retracing the strands. Associativity is again free, since the rule is composition.

\sfigs{0.86}{Closure in $S_4$: after $\sigma$, then after $\tau$, the net effect is again a single
rearrangement.}{\figXt}

\sfigs{0.86}{Identity and inverses in $S_4$ in one picture: doing nothing after $\sigma$ changes no row,
and applying $\sigma\inv$ instead brings every token home.}{\figXu}

\subsubsection*{Cycle notation}

Writing a permutation as a list of images is correct but unreadable. The standard notation records
the \emph{loops} instead.

\begin{defi}
$(i_1\,i_2\,\dots\,i_k)$ denotes the permutation sending $i_1$ to $i_2$, $i_2$ to $i_3$, and so on
round to $i_k$, which goes back to $i_1$; every number not listed is left alone. Such a permutation
is a \defn{$k$-cycle}, and a $2$-cycle is called a \defn{transposition} or a \defn{swap}.
Permutations written side by side, as in $(1\,2\,4)(3\,5)$, are to be multiplied.
\end{defi}

\sfigs{0.86}{The permutation $(1\,2\,4)(3\,5)$ drawn as arrows on the six numbers. Each cycle is a closed
loop, disjoint cycles are loops that share no points, and $6$ is left alone.}{\figXw}

\begin{pitfall}
The same cycle has several names: $(1\,2\,4)$, $(2\,4\,1)$ and $(4\,1\,2)$ are the same permutation,
because the notation lists a loop starting from a chosen point and going round once. The starting
point is a choice; the loop is not.
\end{pitfall}

\sfigs{0.86}{One loop, three names.}{\figXx}

\subsubsection*{Multiplying permutations}

There is one reliable method: push each number through both factors, remembering that the
right-hand factor acts first. Take $(2\,6)\cdot(1\,2\,4)(3\,5)$.

\begin{center}
\small
\begin{tabular}{@{}cccc@{}}
\toprule
number & after $(1\,2\,4)(3\,5)$ & then after $(2\,6)$ & so the product sends\\
\midrule
1 & 2 & 6 & $1\to6$\\
2 & 4 & 4 & $2\to4$\\
3 & 5 & 5 & $3\to5$\\
4 & 1 & 1 & $4\to1$\\
5 & 3 & 3 & $5\to3$\\
6 & 6 & 2 & $6\to2$\\
\bottomrule
\end{tabular}
\end{center}

Now read the answer off as loops. Start at $1$ and follow: $1\to6\to2\to4\to1$, which closes up as
the $4$-cycle $(1\,6\,2\,4)$. The smallest unused number is $3$, and $3\to5\to3$ gives $(3\,5)$.
Everything is accounted for, so
\[(2\,6)\cdot(1\,2\,4)(3\,5)=(1\,6\,2\,4)(3\,5).\]

\sfigs{0.86}{The answer as a loop picture.}{\figXz}

The same computation can be watched happening to the tokens. Each row is a complete snapshot of
where everybody is; the strands between rows show each token's move.

\sfigs{0.86}{One factor at a time. Read the last row: token $1$ sits in slot $6$, token $6$ in slot $2$,
token $2$ in slot $4$, token $4$ in slot $1$, and $3,5$ have swapped. That is the right-hand column
of the table above.}{\figXaa}

\begin{why}
The table and the picture agree because they organise the same information differently. The table
follows one \emph{number} at a time through both factors. The picture follows all six at once,
one \emph{factor} at a time. Neither is more correct; the table is better for computing, the picture
is better for seeing.
\end{why}

\begin{pitfall}
Apply $(2\,6)$ first instead and the first line of the table becomes $1\to1\to2$, and the answer
changes to $(1\,2\,4\,6)(3\,5)$, a different element. The order is a convention, but it is not
optional: it is fixed by Definition~\ref{def:order}.
\end{pitfall}

\subsubsection*{Permutations as diagrams, and why swaps generate}

Draw the slots as $n$ points on a top row and $n$ points on a bottom row, and join each token's
start to its finish by a strand. Nudge the picture so that no two crossings happen at the same
height. Reading the crossings from top to bottom expresses the permutation as a succession of swaps
of \emph{neighbouring} numbers, one per crossing.

\sfig{The permutation $(1\,2\,4)(3\,5)$ as a stack of five neighbour swaps. One row per height, one
swap between consecutive rows.}{\figXab}

\sfig{The bottom row of that stack, read as loops, is the permutation we started
with.}{\figXac}

\sfig{Stacking two diagrams multiplies the permutations. The factor performed first goes on top, so
the picture reads downwards in the same order as the product $\,(2\,6)\cdot(1\,2\,4)(3\,5)$ is
computed. Following a strand from the top row to the bottom row gives that number's image under the
product.}{\figstack}

\begin{prop}\label{prop:swapsgen}
Every element of $S_n$ is a product of transpositions; in fact of transpositions of neighbours.
\end{prop}

\begin{proof}
Given a permutation, sort the tokens into place by repeatedly swapping a neighbouring pair that is
out of order, exactly as in the sorting algorithm the picture above performs. The process ends,
because each swap reduces by one the number of pairs that are in the wrong relative order, and that
number is finite. Reading the swaps in the order performed expresses the permutation as their
product.
\end{proof}

\subsubsection*{The extra properties}

$|S_n|=n!$, because the token in slot $1$ has $n$ possible destinations, the next has $n-1$, and so
on. $S_n$ is not abelian once $n\ge3$: for instance $(1\,2)(2\,3)=(1\,2\,3)$ while
$(2\,3)(1\,2)=(1\,3\,2)$. By Proposition~\ref{prop:swapsgen} it is generated by neighbour swaps, a
generating set of size $n-1$.

And by Theorem~\ref{thm:cubeS4}, $S_4$ is the rotation group of the cube. With cycle notation
available, the dictionary at the end of Section~\ref{sec:def} can now be read properly: a quarter
turn about a face axis is a $4$-cycle such as $(d_1\,d_2\,d_4\,d_3)$, a half turn about a face axis
is a pair of swaps such as $(d_1\,d_4)(d_2\,d_3)$, a third of a turn about a diagonal is a $3$-cycle
such as $(d_2\,d_4\,d_3)$, and a half turn about an edge axis is a single swap such as
$(d_1\,d_2)$.

\begin{rem}
The full symmetry group of the cube, reflections included, has $48$ elements and is \emph{not}
$S_4$. It is $S_4$ paired with a two-element group that records whether a motion is a mirror image.
Products of groups are not needed in this chapter, but that is where the extra factor comes from.
\end{rem}

\begin{tryit}
Compute $(1\,2\,3\,4)\cdot(1\,2\,3\,4)$ and $\big((1\,2)(3\,4)\big)\inv$. Then write
$(1\,3)(2\,4)$ as a product of neighbour swaps, and check your answer by composing them.
\end{tryit}
